\documentclass{beamer}
\usepackage{amsfonts,amsmath,oldgerm}
\usetheme{sintef}
\usepackage{xeCJK}
\setCJKmainfont{Noto Serif CJK SC}

\newcommand{\testcolor}[1]{\colorbox{#1}{\textcolor{#1}{test}}~\texttt{#1}}

\usefonttheme[onlymath]{serif}

\titlebackground*{assets/background}

\newcommand{\hrefcol}[2]{\textcolor{cyan}{\href{#1}{#2}}}

\title{迈向快速且可泛化的开放域决策}
\subtitle{Intern-Decision 作为System-1 决策层}
% \course{Master's Degree in Computer Science}
\author{周家纬、孙哲、赵博涵、张帅宇、李春晓、孙铮浩、谢嘉庆}
% \IDnumber{1234567}
\date{2026 年11 月19 日}

\begin{document}
\maketitle

\begin{frame}{CPU 实验目标与研究重点}
\begin{itemize}
    \item \textbf{实验背景:} 受限于算力资源，初期实验基于 CPU 进行快速迭代，旨在以最小化成本验证算法框架的可行性。
    \item \textbf{关键超参数筛选:} 在受限环境下，对模型关键超参数进行高效寻优，确定适合大规模实验的基准配置。
    \item \textbf{Calibration 有效性验证:} 评估校准策略在 System-1 决策层中的表现，验证置信度校准对决策输出的优化效果。
    \item \textbf{决策可信度分析:} 通过有限资源下的小规模评估，查看效果以便于建立决策机制的基准评估体系，量化分析模型在不同配置下的输出稳定性与可信度。
\end{itemize}
\end{frame}

\section{Recall}

\begin{frame}{研究计划概览}
\begin{itemize}
    \item \textbf{实验 1 (E1)：} 复现 Intern Decision 4b，0.8b 的基线表现
    \item \textbf{实验 2 (E2)：} 量“快”值不值，看 Jev 类模型和同参数下LLM的 CoT 回答差距
    \item \textbf{实验 3 (E3)：} 设计置信度机制，看什么时候该求助
    \item \textbf{实验 4 (E4)：} 在jev上微调后，其能力是否可以迁移（以实验室 gpu 任务排班为例），并且验证是否可以容纳多领域内的知识（旧题重测）
\end{itemize}
\end{frame}

\section{Progress}

\begin{frame}[fragile]{实验环境：E1-E4 自动化构建}
  \begin{columns}[T]
    \begin{column}{0.35\textwidth}
      \begin{itemize}
        \item \textbf{框架现状:} E1-E4 实验脚本及自动化工作流已构建完成。
        \item \textbf{核心特性:}
        \begin{itemize}
          \item 支持多后端 (HF/XTuner)。
          \item 离线模式 (Offline) 支持。
          \item 开启 GPU 环境后，可直接执行脚本启动任务。
        \end{itemize}
      \end{itemize}
    \end{column}
    \begin{column}{0.65\textwidth}
      \tiny
      \begin{verbatim}
#!/usr/bin/env bash
set -euo pipefail
E1_ROOT="$(cd -- "$(dirname -- "${BASH_SOURCE[0]}")/.." && pwd)"
E1_BACKEND="${1:-hf}"
if [[ "$E1_BACKEND" != hf && "$E1_BACKEND" != xtuner ]]; then
  exit 2
fi
cd "$E1_ROOT"
E1_PYTHON="$E1_ROOT/.venv-$E1_BACKEND/bin/python"
export HF_HUB_OFFLINE=1 TRANSFORMERS_OFFLINE=1 HF_DATASETS_OFFLINE=1
export OMP_NUM_THREADS="${OMP_NUM_THREADS:-8}"
if [[ "$E1_BACKEND" == xtuner ]]; then
  export PYTHONPATH="$E1_ROOT:$E1_ROOT/upstream/Intern-Decision"
fi
"$E1_PYTHON" -m decision_lab.offline_bundle check --backend "$E1_BACKEND"
exec "$E1_PYTHON" scripts/run_e1_baseline.py \
  --backend "$E1_BACKEND" \
  --output "$E1_OUTPUT"
\end{verbatim}
    \end{column}
  \end{columns}
\end{frame}

\backmatter
\end{document}
